\documentclass{article}
\usepackage[UTF8]{ctex}
\usepackage{amsmath,amssymb,amsfonts}
\usepackage{enumitem}
\usepackage{geometry}
\geometry{a4paper, margin=1in}
% === 错题本排版优化（不参与题目标识） ===
\setlength{\parindent}{0pt}
\setlength{\parskip}{0.35em plus 0.10em minus 0.08em}
\setlength{\emergencystretch}{3em}
\setlist{itemsep=0.25em, topsep=0.40em, parsep=0pt}
\allowdisplaybreaks[2]
% === 错题本排版优化结束 ===
\begin{document}
\section*{暨南大学2024年高等代数真题}
\section*{148．暨南大学2024年高等代数试卷（回忆版）}

\subsection{题目1}
1．描述代数学基本定理与哈密顿－凯莱定理。

\subsection{题目2}
2．讨论 \(m, p, q\) 满足怎么条件时，有 \(x^{2}+m x+1 \mid x^{3}+p x+q\) ．

\subsection{题目3}
3．若 \(A, B\) 为实对称矩阵，\(A+B\) 是否实对称？\(A B\) 呢？

\subsection{题目4}
4．设 \(Y \in \mathbb{R}^{n}\) 长度为 \(1, A\) 为 \(n\) 阶正定矩阵，\(A\) 特征值最大为2024，证明：\(Y^{\prime} A Y \leq 2024\) 。

\subsection{题目5}
5．是否存在线性变换，使得平面上方不变，下方变到上方，即 \((x, y) \mapsto(x,-y)\) ．

\subsection{题目6}
6．设欧氏空间中的向量 \(\alpha_{1}, \alpha_{2}, \alpha_{3}\) 满足

\[
C=\left(\begin{array}{ll}
\left(\alpha_{1}, \alpha_{1}\right) & \left(\alpha_{1}, \alpha_{2}\right) \\
\left(\alpha_{2}, \alpha_{1}\right) & \left(\alpha_{2}, \alpha_{2}\right) \\
\left(\alpha_{3}, \alpha_{1}\right) & \left(\alpha_{3}, \alpha_{2}\right) \\
\left(\alpha_{3}, \alpha_{3}\right)
\end{array}\right)
\]

可逆，证明：\(\alpha_{1}, \alpha_{2}, \alpha_{3}\) 线性无关。

\subsection{题目7}
7．讨论方程组（具体数值忘了）当 \(k\) 复合什么条件时只有零解．

\subsection{题目8}
8．设 \(X=\{(x, 0) \mid x \in \mathbb{R}\}, Y=\{(0, y) \mid y \in \mathbb{R}\}\) ．

\begin{itemize}
  \item[（1）] 是否有 \(\mathbb{R}^{2}=X \oplus Y\) ？
  \item[（2）] 是否有 \(\mathbb{R}^{2}=X \cup Y\) ？
\end{itemize}

\subsection{题目9}
9．设 \(A, B, C, D\) 为 \(n\) 阶方阵，且 \(A C=C A\) ，证明：\(\left|\begin{array}{ll}A & B \\ C & D\end{array}\right|=|A D-C B|\) 。

\subsection{题目10}
10．设 \(A\) 是 \(n\) 阶实矩阵，证明：\(r\left(A^{\prime} A\right)=r(A)\) 。

\subsection{题目11}
11．设 \(A=X X^{\prime}\) ，其中 \(X \in \mathbb{R}^{n}\) 是 \(n\) 维非零列向量。

\begin{itemize}
  \item[（1）] 证明：\(A\) 为半正定矩阵．
  \item[（2）] 证明：\(A\) 有且仅有一个非零特征值．
\end{itemize}
\end{document}
